% GENERATED FILE. Edit canonical lessons, then run npm run generate:books.
% canonical-source-hash: fnv1a64:02357cd9e8ae7fdd
% canonical-lessons: SA-C465-cancalah, SA-C465-udasinah, SA-C465-pramadyati, SA-C465-avadhanam, SA-C465-udyami

\chapter{Steady and Careless}
\label{ch:sa-steady-and-careless}
\hlchaptermodality{465}

\begin{chapteropening}
\textbf{By the end of this chapter:} \emph{I can say whether someone works carefully.}

Complete the last lesson of chapter 465: I can say whether someone works carefully.
\end{chapteropening}

\section[\texorpdfstring{\sk{चञ्चलः} (cañcalaḥ)}{चञ्चलः (cañcalaḥ)}]{\sk{चञ्चलः} (cañcalaḥ) — fickle}

\label{lesson:SA-C465-cancalah}



\begin{quote}
Before the new one: say the Sanskrit for reviles, then the Sanskrit for prattle.
\end{quote}

\subsection*{You'll want to know: \sk{चञ्चलः}}
\textbf{\sk{चञ्चलः}} — \emph{cañcalaḥ} — \textquotedblleft{}fickle\textquotedblright{}.

\begin{quote}
\sk{वानरस्य} \sk{मनः} \sk{चञ्चलं} \sk{भवति।}
\end{quote}

\begin{quote}
\emph{vānarasya manaḥ cañcalaṃ bhavati.}
\end{quote}

\begin{quote}
A monkey's mind is fickle.
\end{quote}

\begin{cousinweb}
\textbf{\sk{चञ्चलः}} (\emph{cañcalaḥ}) is built on the intensive of the root √cal, \textquotedblleft{}to move\textquotedblright{}, the root of \textbf{\sk{चलति}}: moving to and fro; MW gives \textquotedblleft{}unsteady, inconstant\textquotedblright{} (MW).
\end{cousinweb}

\begin{grammarlens}[title={What you've built}]
One more word for describing what you lived and felt.
\end{grammarlens}

\subsection*{Your turn}
\begin{itemize}
\raggedright
  \item \emph{Say it:} \emph{cañcalaḥ}
  \item \emph{Say it:} \emph{cañcalaḥ}, once more
  \item \emph{Say it:} \emph{vānarasya manaḥ cañcalaṃ bhavati.} — A monkey's mind is fickle.
  \item \emph{Recall:} say the Sanskrit for reviles, then the Sanskrit for prattle, then say \emph{cañcalaḥ} again
\end{itemize}

\begin{tcolorbox}[breakable,colback=teal!4,colframe=teal!35!black,title={Before you move on}]
What does \sk{चञ्चलः} mean? (\textquotedblleft{}Fickle\textquotedblright{}.) Say it once more.
\end{tcolorbox}
\section[\texorpdfstring{\sk{उदासीनः} (udāsīnaḥ)}{उदासीनः (udāsīnaḥ)}]{\sk{उदासीनः} (udāsīnaḥ) — indifferent}

\label{lesson:SA-C465-udasinah}



\begin{quote}
Before the new one: say the Sanskrit for prattle, then the Sanskrit for fickle.
\end{quote}

\subsection*{You'll want to know: \sk{उदासीनः}}
\textbf{\sk{उदासीनः}} — \emph{udāsīnaḥ} — \textquotedblleft{}indifferent\textquotedblright{}.

\begin{quote}
\sk{सः} \sk{लाभे} \sk{हानौ} \sk{च} \sk{उदासीनः} \sk{अस्ति।}
\end{quote}

\begin{quote}
\emph{saḥ lābhe hānau ca udāsīnaḥ asti.}
\end{quote}

\begin{quote}
He is indifferent to gain and loss.
\end{quote}

\begin{cousinweb}
\textbf{\sk{उदासीनः}} (\emph{udāsīnaḥ}) is built from \emph{ud}, \textquotedblleft{}up, apart\textquotedblright{}, and \emph{āsīna}, \textquotedblleft{}sitting\textquotedblright{}, from the root √ās of \textbf{\sk{आस्ते}}: sitting apart; MW gives \textquotedblleft{}indifferent, free from affection\textquotedblright{} (MW).
\end{cousinweb}

\begin{grammarlens}[title={What you've built}]
One more word for describing what you lived and felt.
\end{grammarlens}

\subsection*{Your turn}
\begin{itemize}
\raggedright
  \item \emph{Say it:} \emph{udāsīnaḥ}
  \item \emph{Say it:} \emph{udāsīnaḥ}, once more
  \item \emph{Say it:} \emph{saḥ lābhe hānau ca udāsīnaḥ asti.} — He is indifferent to gain and loss.
  \item \emph{Recall:} say the Sanskrit for prattle, then the Sanskrit for fickle, then say \emph{udāsīnaḥ} again
\end{itemize}

\begin{tcolorbox}[breakable,colback=teal!4,colframe=teal!35!black,title={Before you move on}]
What does \sk{उदासीनः} mean? (\textquotedblleft{}Indifferent\textquotedblright{}.) Say it once more.
\end{tcolorbox}
\section[\texorpdfstring{\sk{प्रमाद्यति} (pramādyati)}{प्रमाद्यति (pramādyati)}]{\sk{प्रमाद्यति} (pramādyati) — is careless}

\label{lesson:SA-C465-pramadyati}



\begin{quote}
Before the new one: say the Sanskrit for fickle, then the Sanskrit for indifferent.
\end{quote}

\subsection*{You'll want to know: \sk{प्रमाद्यति}}
\textbf{\sk{प्रमाद्यति}} — \emph{pramādyati} — \textquotedblleft{}is careless\textquotedblright{}.

\begin{quote}
\sk{चालकः} \sk{कदापि} \sk{न} \sk{प्रमाद्येत्।}
\end{quote}

\begin{quote}
\emph{cālakaḥ kadāpi na pramādyet.}
\end{quote}

\begin{quote}
A driver should never be careless.
\end{quote}

Past (he …): \sk{सः} \sk{प्रमत्तवान्}, \emph{saḥ pramattavān}.

\begin{cousinweb}
\textbf{\sk{प्रमाद्यति}} (\emph{pramādyati}) is built from \emph{pra} and the root √mad, \textquotedblleft{}to be drunk, to rejoice\textquotedblright{}; MW gives \textquotedblleft{}to be careless or negligent\textquotedblright{}, and its ppp \emph{pramatta} is \textquotedblleft{}careless, heedless\textquotedblright{} (MW).
\end{cousinweb}

\begin{grammarlens}[title={What you've built}]
One more word for describing what you lived and felt.
\end{grammarlens}

\subsection*{Your turn}
\begin{itemize}
\raggedright
  \item \emph{Say it:} \emph{pramādyati}
  \item \emph{Say it:} \emph{pramādyati}, once more
  \item \emph{Say it:} \emph{cālakaḥ kadāpi na pramādyet.} — A driver should never be careless.
  \item \emph{Recall:} say the Sanskrit for fickle, then the Sanskrit for indifferent, then say \emph{pramādyati} again
\end{itemize}

\begin{tcolorbox}[breakable,colback=teal!4,colframe=teal!35!black,title={Before you move on}]
What does \sk{प्रमाद्यति} mean? (\textquotedblleft{}Is careless\textquotedblright{}.) Say it once more.
\end{tcolorbox}
\section[\texorpdfstring{\sk{अवधानम्} (avadhānam)}{अवधानम् (avadhānam)}]{\sk{अवधानम्} (avadhānam) — attention}

\label{lesson:SA-C465-avadhanam}



\begin{quote}
Before the new one: say the Sanskrit for indifferent, then the Sanskrit for is careless.
\end{quote}

\subsection*{You'll want to know: \sk{अवधानम्}}
\textbf{\sk{अवधानम्}} — \emph{avadhānam} — \textquotedblleft{}attention\textquotedblright{}.

\begin{quote}
\sk{अवधानेन} \sk{पाठं} \sk{शृणोतु।}
\end{quote}

\begin{quote}
\emph{avadhānena pāṭhaṃ śṛṇotu.}
\end{quote}

\begin{quote}
Listen to the lesson with attention.
\end{quote}

A neuter noun.

\begin{cousinweb}
\textbf{\sk{अवधानम्}} (\emph{avadhānam}) is built from \emph{ava}, \textquotedblleft{}down\textquotedblright{}, and the root √dhā, \textquotedblleft{}to put\textquotedblright{}: putting the mind down on a thing; MW gives \textquotedblleft{}attention, attentiveness\textquotedblright{} (MW).
\end{cousinweb}

\begin{grammarlens}[title={What you've built}]
One more word for describing what you lived and felt.
\end{grammarlens}

\subsection*{Your turn}
\begin{itemize}
\raggedright
  \item \emph{Say it:} \emph{avadhānam}
  \item \emph{Say it:} \emph{avadhānam}, once more
  \item \emph{Say it:} \emph{avadhānena pāṭhaṃ śṛṇotu.} — Listen to the lesson with attention.
  \item \emph{Recall:} say the Sanskrit for indifferent, then the Sanskrit for is careless, then say \emph{avadhānam} again
\end{itemize}

\begin{tcolorbox}[breakable,colback=teal!4,colframe=teal!35!black,title={Before you move on}]
What does \sk{अवधानम्} mean? (\textquotedblleft{}Attention\textquotedblright{}.) Say it once more.
\end{tcolorbox}
\section[\texorpdfstring{\sk{उद्यमी} (udyamī)}{उद्यमी (udyamī)}]{\sk{उद्यमी} (udyamī) — industrious}

\label{lesson:SA-C465-udyami}



\begin{quote}
Before the new one: say the Sanskrit for is careless, then the Sanskrit for attention.
\end{quote}

\subsection*{You'll want to know: \sk{उद्यमी}}
\textbf{\sk{उद्यमी}} — \emph{udyamī} — \textquotedblleft{}industrious\textquotedblright{}.

\begin{quote}
\sk{उद्यमी} \sk{कृषकः} \sk{प्रातः} \sk{क्षेत्रं} \sk{गच्छति।}
\end{quote}

\begin{quote}
\emph{udyamī kṛṣakaḥ prātaḥ kṣetraṃ gacchati.}
\end{quote}

\begin{quote}
The industrious farmer goes to the field early.
\end{quote}

\begin{cousinweb}
\textbf{\sk{उद्यमी}} (\emph{udyamī}) is built on \emph{udyama}, \textquotedblleft{}exertion\textquotedblright{}, from \emph{ud} and the root √yam, \textquotedblleft{}to hold\textquotedblright{}; MW gives \textquotedblleft{}persevering, making effort, active\textquotedblright{} (MW).
\end{cousinweb}

\begin{grammarlens}[title={What you've built}]
One more word for describing what you lived and felt.
\end{grammarlens}

\subsection*{Your turn}
\begin{itemize}
\raggedright
  \item \emph{Say it:} \emph{udyamī}
  \item \emph{Say it:} \emph{udyamī}, once more
  \item \emph{Say it:} \emph{udyamī kṛṣakaḥ prātaḥ kṣetraṃ gacchati.} — The industrious farmer goes to the field early.
  \item \emph{Recall:} say the Sanskrit for is careless, then the Sanskrit for attention, then say \emph{udyamī} again
\end{itemize}

\begin{tcolorbox}[breakable,colback=teal!4,colframe=teal!35!black,title={Before you move on}]
What does \sk{उद्यमी} mean? (\textquotedblleft{}Industrious\textquotedblright{}.) Say it once more.
\end{tcolorbox}
